% Options for packages loaded elsewhere
\PassOptionsToPackage{unicode}{hyperref}
\PassOptionsToPackage{hyphens}{url}
%
\documentclass[
]{article}
\usepackage{amsmath,amssymb}
\usepackage{iftex}
\ifPDFTeX
  \usepackage[T1]{fontenc}
  \usepackage[utf8]{inputenc}
  \usepackage{textcomp} % provide euro and other symbols
\else % if luatex or xetex
  \usepackage{unicode-math} % this also loads fontspec
  \defaultfontfeatures{Scale=MatchLowercase}
  \defaultfontfeatures[\rmfamily]{Ligatures=TeX,Scale=1}
\fi
\usepackage{lmodern}
\ifPDFTeX\else
  % xetex/luatex font selection
\fi
% Use upquote if available, for straight quotes in verbatim environments
\IfFileExists{upquote.sty}{\usepackage{upquote}}{}
\IfFileExists{microtype.sty}{% use microtype if available
  \usepackage[]{microtype}
  \UseMicrotypeSet[protrusion]{basicmath} % disable protrusion for tt fonts
}{}
\makeatletter
\@ifundefined{KOMAClassName}{% if non-KOMA class
  \IfFileExists{parskip.sty}{%
    \usepackage{parskip}
  }{% else
    \setlength{\parindent}{0pt}
    \setlength{\parskip}{6pt plus 2pt minus 1pt}}
}{% if KOMA class
  \KOMAoptions{parskip=half}}
\makeatother
\usepackage{xcolor}
\usepackage[margin=1in]{geometry}
\usepackage{color}
\usepackage{fancyvrb}
\newcommand{\VerbBar}{|}
\newcommand{\VERB}{\Verb[commandchars=\\\{\}]}
\DefineVerbatimEnvironment{Highlighting}{Verbatim}{commandchars=\\\{\}}
% Add ',fontsize=\small' for more characters per line
\usepackage{framed}
\definecolor{shadecolor}{RGB}{248,248,248}
\newenvironment{Shaded}{\begin{snugshade}}{\end{snugshade}}
\newcommand{\AlertTok}[1]{\textcolor[rgb]{0.94,0.16,0.16}{#1}}
\newcommand{\AnnotationTok}[1]{\textcolor[rgb]{0.56,0.35,0.01}{\textbf{\textit{#1}}}}
\newcommand{\AttributeTok}[1]{\textcolor[rgb]{0.13,0.29,0.53}{#1}}
\newcommand{\BaseNTok}[1]{\textcolor[rgb]{0.00,0.00,0.81}{#1}}
\newcommand{\BuiltInTok}[1]{#1}
\newcommand{\CharTok}[1]{\textcolor[rgb]{0.31,0.60,0.02}{#1}}
\newcommand{\CommentTok}[1]{\textcolor[rgb]{0.56,0.35,0.01}{\textit{#1}}}
\newcommand{\CommentVarTok}[1]{\textcolor[rgb]{0.56,0.35,0.01}{\textbf{\textit{#1}}}}
\newcommand{\ConstantTok}[1]{\textcolor[rgb]{0.56,0.35,0.01}{#1}}
\newcommand{\ControlFlowTok}[1]{\textcolor[rgb]{0.13,0.29,0.53}{\textbf{#1}}}
\newcommand{\DataTypeTok}[1]{\textcolor[rgb]{0.13,0.29,0.53}{#1}}
\newcommand{\DecValTok}[1]{\textcolor[rgb]{0.00,0.00,0.81}{#1}}
\newcommand{\DocumentationTok}[1]{\textcolor[rgb]{0.56,0.35,0.01}{\textbf{\textit{#1}}}}
\newcommand{\ErrorTok}[1]{\textcolor[rgb]{0.64,0.00,0.00}{\textbf{#1}}}
\newcommand{\ExtensionTok}[1]{#1}
\newcommand{\FloatTok}[1]{\textcolor[rgb]{0.00,0.00,0.81}{#1}}
\newcommand{\FunctionTok}[1]{\textcolor[rgb]{0.13,0.29,0.53}{\textbf{#1}}}
\newcommand{\ImportTok}[1]{#1}
\newcommand{\InformationTok}[1]{\textcolor[rgb]{0.56,0.35,0.01}{\textbf{\textit{#1}}}}
\newcommand{\KeywordTok}[1]{\textcolor[rgb]{0.13,0.29,0.53}{\textbf{#1}}}
\newcommand{\NormalTok}[1]{#1}
\newcommand{\OperatorTok}[1]{\textcolor[rgb]{0.81,0.36,0.00}{\textbf{#1}}}
\newcommand{\OtherTok}[1]{\textcolor[rgb]{0.56,0.35,0.01}{#1}}
\newcommand{\PreprocessorTok}[1]{\textcolor[rgb]{0.56,0.35,0.01}{\textit{#1}}}
\newcommand{\RegionMarkerTok}[1]{#1}
\newcommand{\SpecialCharTok}[1]{\textcolor[rgb]{0.81,0.36,0.00}{\textbf{#1}}}
\newcommand{\SpecialStringTok}[1]{\textcolor[rgb]{0.31,0.60,0.02}{#1}}
\newcommand{\StringTok}[1]{\textcolor[rgb]{0.31,0.60,0.02}{#1}}
\newcommand{\VariableTok}[1]{\textcolor[rgb]{0.00,0.00,0.00}{#1}}
\newcommand{\VerbatimStringTok}[1]{\textcolor[rgb]{0.31,0.60,0.02}{#1}}
\newcommand{\WarningTok}[1]{\textcolor[rgb]{0.56,0.35,0.01}{\textbf{\textit{#1}}}}
\usepackage{longtable,booktabs,array}
\usepackage{calc} % for calculating minipage widths
% Correct order of tables after \paragraph or \subparagraph
\usepackage{etoolbox}
\makeatletter
\patchcmd\longtable{\par}{\if@noskipsec\mbox{}\fi\par}{}{}
\makeatother
% Allow footnotes in longtable head/foot
\IfFileExists{footnotehyper.sty}{\usepackage{footnotehyper}}{\usepackage{footnote}}
\makesavenoteenv{longtable}
\usepackage{graphicx}
\makeatletter
\def\maxwidth{\ifdim\Gin@nat@width>\linewidth\linewidth\else\Gin@nat@width\fi}
\def\maxheight{\ifdim\Gin@nat@height>\textheight\textheight\else\Gin@nat@height\fi}
\makeatother
% Scale images if necessary, so that they will not overflow the page
% margins by default, and it is still possible to overwrite the defaults
% using explicit options in \includegraphics[width, height, ...]{}
\setkeys{Gin}{width=\maxwidth,height=\maxheight,keepaspectratio}
% Set default figure placement to htbp
\makeatletter
\def\fps@figure{htbp}
\makeatother
\setlength{\emergencystretch}{3em} % prevent overfull lines
\providecommand{\tightlist}{%
  \setlength{\itemsep}{0pt}\setlength{\parskip}{0pt}}
\setcounter{secnumdepth}{5}
\ifLuaTeX
  \usepackage{selnolig}  % disable illegal ligatures
\fi
\usepackage{bookmark}
\IfFileExists{xurl.sty}{\usepackage{xurl}}{} % add URL line breaks if available
\urlstyle{same}
\hypersetup{
  pdftitle={Étude de la distribution des céphalophes dans le bassin du Congo},
  pdfauthor={Nom de l'équipe},
  hidelinks,
  pdfcreator={LaTeX via pandoc}}

\title{Étude de la distribution des céphalophes dans le bassin du Congo}
\author{Nom de l'équipe}
\date{2025-04-16}

\begin{document}
\maketitle

{
\setcounter{tocdepth}{2}
\tableofcontents
}
\section{Informations légales}\label{informations-luxe9gales}

Ce document présente les résultats de l'étude menée sur l'abondance et la distribution des espèces de céphalophes dans cinq villages du bassin du Congo. Il inclut le traitement des données issues de pièges photographiques, l'analyse de la pression anthropique et des facteurs environnementaux. Cette analyse a été conduite dans le cadre de l'évaluation du temps 1 à GEC.

\section{Résumé de l'étude}\label{ruxe9sumuxe9-de-luxe9tude}

Notre étude explore la distribution des céphalophes dans le bassin du Congo en lien avec des facteurs environnementaux et anthropiques. À partir d'observations issues de pièges photographiques, nous analysons l'effet du relief, de la végétation, des routes et de la proximité des villages sur l'abondance des céphalophes. Nos résultats révèlent une influence significative de la pression humaine, modélisée par une analyse en composantes principales, sur leur distribution. Nous proposons des pistes pour la conservation et la gestion durable de cette espèce.

\section{Introduction}\label{introduction}

\subsection{Contexte}\label{contexte}

\subsection{Problématique}\label{probluxe9matique}

\subsection{Objectifs}\label{objectifs}

\subsection{Hypothèses}\label{hypothuxe8ses}

\begin{enumerate}
\def\labelenumi{\arabic{enumi}.}
\tightlist
\item
  La pression anthropique réduit la présence des céphalophes
\item
  La couverture végétale dense est favorable à leur présence
\item
  Le relief influence l'accessibilité pour la faune
\end{enumerate}

\section{Méthodologie}\label{muxe9thodologie}

\subsection{Données brutes utilisées}\label{donnuxe9es-brutes-utilisuxe9es}

\subsection{Préparation des données}\label{pruxe9paration-des-donnuxe9es}

\subsection{Modélisation}\label{moduxe9lisation}

\subsection{Limites}\label{limites}

\begin{itemize}
\tightlist
\item
  Données manquantes
\item
  Saison non prise en compte
\end{itemize}

\section{Résultats de l'analyse}\label{ruxe9sultats-de-lanalyse}

\subsection{Géographie du territoire}\label{guxe9ographie-du-territoire}

\begin{itemize}
\tightlist
\item
  Carte des villages et stations (à insérer)
\end{itemize}

Présentation des 5 villages : {[}à compléter{]}

\begin{itemize}
\tightlist
\item
  Localisation dans le bassin du Congo
\item
  Couverture végétale MODIS
\item
  Répartition des stations
\end{itemize}

\subsection{Abondance des céphalophes}\label{abondance-des-cuxe9phalophes}

\begin{itemize}
\tightlist
\item
  Histogrammes d'abondance par village
\item
  Heatmaps des observations par station
\end{itemize}

\subsection{Effets environnementaux et anthropiques}\label{effets-environnementaux-et-anthropiques}

\begin{itemize}
\item
  Régression multiple
\item
  Corrélogrammes
\item
  Nuages de points avec axes ACP
\item
  Modèle linéaire sur la densité en fonction du relief et végétation
\item
  GLM avec pression anthropique en prédicteur
\end{itemize}

\begin{Shaded}
\begin{Highlighting}[]
\CommentTok{\# Exemple de modèle}
\CommentTok{\# glm(abondance \textasciitilde{} pression\_anthropique + relief + vegetation, family=poisson)}
\end{Highlighting}
\end{Shaded}

\section{Discussion}\label{discussion}

\begin{itemize}
\tightlist
\item
  Validité des hypothèses
\end{itemize}

L'hypothèse 1 est confirmée : proximité des camps réduit l'abondance.
L'hypothèse 2 est partiellement validée.
L'heure d'observation influe peu, sauf tôt le matin.

\begin{itemize}
\tightlist
\item
  Comparaison inter-village
\item
  Études similaires
\item
  Recommandations de gestion
\end{itemize}

\section{Conclusion}\label{conclusion}

\begin{itemize}
\tightlist
\item
  Synthèse des résultats
\item
  Recommandations pratiques
\item
  Perspectives futures
\end{itemize}

\section{Annexes}\label{annexes}

\begin{itemize}
\tightlist
\item
  Résultats ACP complets
\item
  Script R (en partie ou en annexe)
\end{itemize}

\section{L'équipe}\label{luxe9quipe}

\begin{longtable}[]{@{}lll@{}}
\toprule\noalign{}
Nom & Rôle & Contribution \\
\midrule\noalign{}
\endhead
\bottomrule\noalign{}
\endlastfoot
\ldots{} & \ldots{} & \ldots{} \\
\end{longtable}

\section{Bibliographie}\label{bibliographie}

\begin{quote}
À compléter manuellement (Zotero ou BibTeX)
\end{quote}

\#Résumé

\#\#Contexte

La biodiversité des forêts tropicales africaines est menacée par l'expansion humaine. Les céphalophes sont de bons indicateurs de l'état des écosystèmes.

\#\#Objectif

Évaluer l'influence des pressions humaines et des caractéristiques environnementales sur la présence des céphalophes.

\#\#Méthodes

Analyse statistique des données collectées dans cinq villages, transformation des données temporelles et création d'un indice de pression anthropique via ACP.

\#\#Résultats clés

Les villages diffèrent significativement en termes de densité de céphalophes. L'abondance est influencée à la fois par les variables environnementales (relief, végétation) et les pressions humaines.

\#\#Conclusion principale : Les pressions anthropiques affectent la présence des céphalophes, en interaction avec les conditions environnementales.

\#L'introduction

\subsection{Contexte de l'étude}\label{contexte-de-luxe9tude}

Le bassin du Congo constitue la deuxième plus grande forêt tropicale du monde. Il abrite une faune exceptionnelle, mais fortement menacée par la pression humaine croissante.

\subsection{Problématique}\label{probluxe9matique-1}

Comment les pressions anthropiques et les variables environnementales influencent-elles la répartition des céphalophes ?

\subsection{Objectifs}\label{objectifs-1}

\begin{itemize}
\item
  Comprendre les facteurs influençant l'abondance des céphalophes
\item
  Identifier les variables environnementales et anthropiques pertinentes
\item
  \textbf{Général} : Comprendre l'influence combinée de l'environnement et de l'homme sur les céphalophes.
\item
  \textbf{Spécifiques} :

  \begin{itemize}
  \tightlist
  \item
    Déterminer les facteurs expliquant l'abondance locale.
  \item
    Identifier les zones à forte pression humaine.
  \item
    Comparer les profils écologiques entre villages.
  \end{itemize}
\end{itemize}

\subsection{Hypothèses}\label{hypothuxe8ses-1}

\begin{enumerate}
\def\labelenumi{\arabic{enumi}.}
\tightlist
\item
  Les villages les plus proches des routes ou camps ont une moindre densité de céphalophes.
\item
  La végétation dense favorise leur présence.
\item
  L'heure de la journée influence l'activité observée.
\end{enumerate}

\section{Méthodologie}\label{muxe9thodologie-1}

\subsection{Données brutes utilisées}\label{donnuxe9es-brutes-utilisuxe9es-1}

Deux jeux de données ont été extraits :
- \textbf{data\_especes} : observations de céphalophes (heure, station, espèce)
- \textbf{data\_stations} : variables environnementales et distances (route, village, camp, eau)
- \texttt{data\_especes} : observations issues de pièges photo
- \texttt{data\_stations} : variables environnementales et anthropiques

\subsection{Préparation des données}\label{pruxe9paration-des-donnuxe9es-1}

\begin{itemize}
\tightlist
\item
  Traitement des heures avec \texttt{lubridate}
\item
  Correction des noms de station pour jointure correcte
\item
  Construction de variables composites (ACP)
\end{itemize}

\subsection{Transformations}\label{transformations}

\begin{itemize}
\tightlist
\item
  Transformation des heures en format \texttt{hms}
\item
  Création d'un score composite de pression anthropique (ACP sur les proximités)
\end{itemize}

\subsection{Méthodes}\label{muxe9thodes}

\begin{itemize}
\item
  ACP sur les distances
\item
  Régression linéaire (abondance \textasciitilde{} axe1 + altitude)
\item
  Modèles GLM et mixtes
\item
  Analyse d'abondance : GLM, LMER
\item
  Analyse multivariée : ACP, regroupements
\end{itemize}

\subsection{Limites}\label{limites-1}

\begin{itemize}
\tightlist
\item
  Données manquantes sur certaines stations
\item
  Effets potentiels de la météo ou des saisons non inclus
\end{itemize}

\section{Résultats de l'analyse}\label{ruxe9sultats-de-lanalyse-1}

\section{Discussions}\label{discussions}

\begin{itemize}
\tightlist
\item
  L'hypothèse 1 est confirmée : proximité des camps réduit l'abondance.
\item
  L'hypothèse 2 est partiellement validée.
\item
  L'heure d'observation influe peu, sauf tôt le matin.
\end{itemize}

Les céphalophes sont sensibles à l'impact humain. Les résultats appellent à une meilleure gestion de l'espace rural et forestier. L'outil d'analyse développé pourrait être déployé ailleurs.

\end{document}
